\chapter{La carte nord-américaine}\label{ch:nw:the-north-american-map}

Trois centres de données dans le New Jersey, à quelques dizaines de kilomètres les uns des autres, hébergent les moteurs d'appariement du groupe NYSE, de Nasdaq et des bourses américaines
d'actions et d'options de Cboe~; un quatrième, à Aurora, à l'ouest de Chicago, héberge le Globex de CME. La lumière dans le vide met \qty{3.94}{\milli\second}
pour aller d'Aurora au bâtiment de Nasdaq à Carteret, la lumière dans le verre \qty{5.76}{\milli\second}. En 2010, la fibre la plus rapide entre
Chicago et Carteret affichait \qty{6.65}{\milli\second}~; en 2016, un réseau de liaisons hertziennes micro-ondes acheminait les données d'Aurora à Carteret en
\qty{3.982}{\milli\second}, soit 1.1~\% au-dessus du temps que met la lumière dans l'air le long du plus court chemin à la surface de la Terre. L'histoire des réseaux de
négociation nord-américains est celle de la réduction de cet écart, et ce chapitre la mesure.

Le chapitre~9 a décrit ce que vend le bâtiment d'une plateforme. Ce chapitre place ces bâtiments sur une carte~: quelle plateforme apparie où, à quelle distance se trouvent les
bâtiments, combien de temps la lumière met entre eux dans le vide, dans l'air et dans le verre, et à quel point les routes publiées s'approchent de ces planchers.
Chaque coordonnée est une ligne du registre avec sa source, chaque distance est calculée sur l'ellipsoïde WGS-84, et chaque carte est tracée
à partir de coordonnées calculées. Les chapitres~11 et~12 font de même pour l'Europe et l'Asie, et les chapitres~13 et~14 regardent à l'intérieur des routes.

\section{Le triangle du New Jersey}

\begin{definition}[Distance géodésique]\label{def:nw:the-north-american-map:geodesic}
La \emph{distance géodésique}\index{distance géodésique} entre deux points de la Terre est la longueur du plus court chemin qui les relie le long
de la surface d'un ellipsoïde de référence (pour les cartes et le positionnement par satellite, le WGS-84), calculée à partir de leurs latitudes et de leurs longitudes.
\end{definition}

La géodésique est le chemin que suivraient un câble ou une ligne de pylônes si rien ne faisait obstacle~: ni rivières, ni terrains privés, ni collines, ni
autorisations. C'est l'étalon naturel d'une route de négociation, et ce n'est pas le plus court chemin à travers l'espace~: la corde droite à travers la
Terre entre Aurora et Carteret est plus courte de \qty{1.68}{\kilo\metre}, soit \qty{5.6}{\micro\second} de lumière, et personne ne va la creuser.

\begin{dated}{2026-09}{Qui apparie où en Amérique du Nord}\label{dat:nw:the-north-american-map:sites}
\begin{itemize}
\item \textbf{Mahwah, New Jersey} (1700 McArthur Boulevard)~: le U.S. Liquidity Center d'ICE, qui héberge les marchés du groupe NYSE et les
  flux consolidés des plans CTA et OPRA.
\item \textbf{Carteret, New Jersey} (NY11, 1400 Federal Boulevard, un bâtiment Equinix)~: le centre de données de Nasdaq, en extension (NY11-4) pour la
  puissance électrique et les baies.
\item \textbf{Secaucus, New Jersey} (Equinix NY5, 800 Secaucus Road)~: les bourses américaines d'actions et d'options de Cboe~; Equinix NY4, au 755
  Secaucus Road, à quelques centaines de mètres, est l'un des points vers lesquels ICE livre les données de marché de Toronto.
\item \textbf{Aurora, Illinois} (2905 Diehl Road)~: le plus grand centre de données de CME Group, loué à CyrusOne jusqu'en 2031~; CME et Google Cloud
  construisent à Aurora une région cloud privée et un site de \omterm{def:nw:colocation-products-and-how-they-are-sold:colo}{colocation}.
\item \textbf{Chicago} (350 Cermak)~: le centre de données secondaire de Cboe pour toutes ses plateformes américaines.
\item \textbf{Markham, Ontario}~: le centre de données principal de TMX (TSX, TSX Venture, TSX Alpha, Montreal Exchange), exploité par TMX~; son adresse postale
  n'a pu être confirmée à partir d'une source primaire et est signalée comme approximative.
\end{itemize}
\end{dated}

La \Cref{fig:nw:the-north-american-map:nj} représente le triangle, tracé à partir des coordonnées de la table des sites. Mahwah est l'angle nord, à 55.6
kilomètres de Carteret et 34.5 de Secaucus~; Secaucus et Carteret sont distants de 25.9 kilomètres. NY4 et NY5 se trouvent à 320 mètres l'un de
l'autre, soit un seul point à cette échelle. Dans le vide, les côtés du triangle prennent 185, 115 et \qty{86}{\micro\second}~; dans une fibre droite, la moitié en plus.

\begin{omfigure}
\begin{tikzpicture}
\begin{axis}[axis equal image,width=10cm,xmin=-30,xmax=58,ymin=-32,ymax=38,xlabel={km à l'est de 74.15\textdegree{} W},
  ylabel={km au nord de 40.8\textdegree{} N},tick label style={font=\footnotesize},label style={font=\small},grid=major,grid style={gray!20},
  table/col sep=comma,unbounded coords=jump]
\addplot[omDef,thick,no markers] table[x=x,y=y]{figdata/networks/10-the-north-american-map/nj_links.csv};
\addplot[only marks,mark=*,mark size=2.6pt,omThm,visualization depends on={value \thisrow{anchor}\as\anc},nodes near coords,
  point meta=explicit symbolic,every node near coord/.append style={font=\footnotesize,anchor=\anc,black}]
  table[x=x,y=y,meta=name,col sep=comma,row sep=\\]{id,x,y,name,anchor\\ mahwah,-0.273,30.981,Mahwah (NYSE),south\\ ny4,6.760,-2.606,Secaucus~: NY4 et NY5 (Cboe),west\\ carteret,-7.865,-23.828,Carteret (Nasdaq),north\\};
\addplot[only marks,mark=none,nodes near coords,point meta=explicit symbolic,
  every node near coord/.append style={font=\scriptsize,fill=white,inner sep=1pt,anchor=center,text=omMeth}]
  table[x=x,y=y,meta=text]{figdata/networks/10-the-north-american-map/nj_labels.csv};
\end{axis}
\end{tikzpicture}

\omcaption{Le triangle du New Jersey à partir de coordonnées calculées~: chaque site à partir de son adresse postale sourcée et d'un géocodeur cité
(\cref{dat:nw:the-north-american-map:sites}), chaque côté étiqueté de sa \omterm{def:nw:the-north-american-map:geodesic}{distance géodésique} WGS-84 et de son temps de lumière dans le vide. La projection est
une projection équirectangulaire locale centrée sur 40.8\textdegree{}~N, 74.15\textdegree{}~W. Données~: \texttt{fig\_map.py} sur
\texttt{firm.geomap}.}\label{fig:nw:the-north-american-map:nj}
\end{omfigure}

Le triangle fixe l'arithmétique du trading d'actions et d'options américaines. Le flux consolidé du plan CTA est produit à Mahwah, donc un
changement de cotation chez Nasdaq sur une des actions de ce plan doit aller de Carteret à Mahwah, y être traité, et revenir avant qu'une firme installée à
Carteret le voie dans le flux consolidé~: au moins deux fois \qty{185.4}{\micro\second}, soit \qty{371}{\micro\second} de lumière dans le vide, avant tout traitement. La même firme lit
le flux direct de Nasdaq dans son propre bâtiment. L'écart entre les flux directs et le processeur d'information sur les titres, que décrit le
livre~1, tient en partie à la géographie, et aucun effort d'ingénierie sur le SIP n'y change la part géographique.

Un routeur à Carteret qui envoie un morceau à Cboe, à Secaucus, et un
autre au NYSE à Mahwah atteint Secaucus en premier~: \qty{126.6}{\micro\second} de fibre droite contre \qty{271}{\micro\second}. Le routage synchronisé
retarde le morceau destiné à Secaucus de la différence, ici environ \qty{144}{\micro\second} le long des géodésiques, en pratique de la différence entre les latences mesurées
des routes réelles de la firme, qui sont plus longues. Chaque course à la latence entre les trois bâtiments (l'arbitrage de latence, au sens du
livre~11), se joue sur ces distances.

\begin{table}[ht]
\centering\small
\begin{tabular}{llrrrr}
\toprule
De & À & Géodésique & Vide & Air & Fibre \\
 & & (km) & (\si{\micro\second}) & (\si{\micro\second}) & (\si{\micro\second}) \\
\midrule
Mahwah & Carteret & 55.6 & 185.4 & 185.4 & 271.0 \\
Mahwah & Secaucus NY4 & 34.5 & 114.9 & 115.0 & 168.0 \\
Secaucus NY4 & Carteret & 25.9 & 86.3 & 86.3 & 126.2 \\
Aurora & Carteret & 1\,180.9 & 3\,939.0 & 3\,940.1 & 5\,758.8 \\
Aurora & Mahwah & 1\,179.0 & 3\,932.6 & 3\,933.8 & 5\,749.4 \\
Aurora & Secaucus NY5 & 1\,191.1 & 3\,973.1 & 3\,974.3 & 5\,808.7 \\
Aurora & 350 Cermak & 52.3 & 174.5 & 174.5 & 255.1 \\
350 Cermak & Carteret & 1\,129.2 & 3\,766.7 & 3\,767.8 & 5\,506.9 \\
Markham & Mahwah & 523.7 & 1\,747.0 & 1\,747.6 & 2\,554.2 \\
Markham & Secaucus NY4 & 550.2 & 1\,835.3 & 1\,835.9 & 2\,683.3 \\
Markham & Carteret & 553.1 & 1\,844.9 & 1\,845.4 & 2\,697.2 \\
\bottomrule
\end{tabular}
\caption{Distances et \omterm{def:nw:the-north-american-map:floor}{planchers de latence} à sens unique entre les sites de l'\cref{dat:nw:the-north-american-map:sites}~: géodésiques WGS-84 (méthode
inverse de Vincenty) et le temps que met la lumière le long de celles-ci dans le vide, dans l'air ($n = 1.0003$) et dans une fibre monomode standard ($n_g = 1.4620$).
Données~: \texttt{nw\_map.pair\_rows()}.}\label{tab:nw:the-north-american-map:pairs}
\end{table}

\section{Le site des contrats à terme de la région de Chicago}

Le 10-K de CME Group désigne le bâtiment d'Aurora comme son plus grand centre de données, loué à CyrusOne dans le cadre d'un bail qui expire en 2031, et décrit une
région cloud privée Google Cloud et un site de \omterm{def:nw:colocation-products-and-how-they-are-sold:colo}{colocation} que CME et Google construisent à Aurora pour héberger ses marchés~; le chapitre~16 revient sur
ce qu'un moteur d'appariement dans une région cloud change. Aurora n'a pas toujours été le site. Laughlin, Aguirre et Grundfest, en lisant les
horodatages propres aux bourses, situent le transfert du moteur d'appariement de CME depuis 350 Cermak, en ville, vers Aurora en 2010, et le transfert de la
\omterm{def:nw:colocation-products-and-how-they-are-sold:colo}{colocation} vers Aurora le 29 janvier 2012~; jusque-là, une firme colocalisée à Cermak, avec une liaison hertzienne micro-ondes partant d'Aurora, payait un
aller-retour métropolitain entre les deux. La distance qui les sépare est de \qty{52.3}{\kilo\metre}, soit \qty{255}{\micro\second} de fibre droite dans chaque sens, si bien que l'aller-retour
en fibre durait au moins une demi-milliseconde, le chiffre que leur article avance.

Du côté de Chicago, on trouve donc un bâtiment pour les contrats à terme et un second site en ville, qui reste le centre de données secondaire de Cboe pour toutes ses
plateformes américaines et un point de présence pour les réseaux. La course la plus suivie du continent relie ce bâtiment unique à trois autres sites~: les prix des contrats à terme sur indices
actions formés à Chicago devancent les prix au comptant formés dans le New Jersey, comme le rappelle le même article, et chaque firme qui négocie sur cette avance
a besoin du chemin le plus rapide d'Aurora vers chacun des sommets du triangle.

\begin{omfigure}
\begin{tikzpicture}
\begin{axis}[axis equal image,width=12.5cm,xmin=-850,xmax=1120,ymin=-230,ymax=320,xlabel={km à l'est de 81\textdegree{} O},
  ylabel={km au nord de 41.8\textdegree{} N},tick label style={font=\footnotesize,/pgf/number format/1000 sep={\,}},label style={font=\small},grid=major,grid style={gray!20},
  table/col sep=comma,unbounded coords=jump]
\addplot[omDef,thick,no markers] table[x=x,y=y]{figdata/networks/10-the-north-american-map/na_links.csv};
\addplot[only marks,mark=*,mark size=2.4pt,omThm,visualization depends on={value \thisrow{anchor}\as\anc},nodes near coords,
  point meta=explicit symbolic,every node near coord/.append style={font=\footnotesize,anchor=\anc,black}]
  table[x=x,y=y,meta=name,col sep=comma,row sep=\\]{id,x,y,name,anchor\\ aurora,-601.116,-0.216,Aurora (CME),north\\ cermak,-549.235,5.944,350 Cermak,south\\ markham,137.811,223.106,Markham (TMX),south\\ mahwah,568.188,-79.593,Mahwah (NYSE),west\\ carteret,560.711,-134.402,Carteret (Nasdaq),west\\};
\addplot[only marks,mark=none,nodes near coords,point meta=explicit symbolic,
  every node near coord/.append style={font=\scriptsize,fill=white,inner sep=1pt,anchor=center,text=omMeth}]
  table[x=x,y=y,meta=text]{figdata/networks/10-the-north-american-map/na_labels.csv};
\end{axis}
\end{tikzpicture}

\omcaption{Chicago, Toronto et New Jersey à partir de coordonnées calculées, chaque liaison étant étiquetée par sa \omterm{def:nw:the-north-american-map:geodesic}{distance géodésique} et son temps de propagation de la lumière dans le vide.
À cette échelle, le triangle du New Jersey de la \cref{fig:nw:the-north-american-map:nj} se réduit à deux points (Mahwah au-dessus, Carteret en dessous). Une
projection équirectangulaire locale centrée sur 41.8\textdegree{}~N, 81\textdegree{}~O~; sur 1\,500~km, elle déforme les formes de quelques pour cent, une erreur acceptable pour un schéma
mais pas pour un calcul de distance (les distances sont calculées sur l'ellipsoïde). Données~: \texttt{fig\_map.py} sur
\texttt{firm.geomap}.}\label{fig:nw:the-north-american-map:na}
\end{omfigure}

\section{Toronto et le reste du continent}

TMX exploite son propre centre de données à Markham, au nord de Toronto~: ses pages décrivent un bâtiment périurbain exploité pour des clients des marchés de capitaux,
neutre vis-à-vis des opérateurs et des fournisseurs, avec des services métropolitains en ondes millimétriques et des liaisons hertziennes micro-ondes transfrontalières sur place. La page d'ICE évoque les actions cotées sur les deux places
et les fonds indiciels cotés (ETF) qui se négocient à Toronto comme à New York, et ICE vend un service sans fil entre les deux~: données de la TSX livrées à Mahwah, à Secaucus NY4 et à Carteret, données américaines livrées à Markham, avec une bande passante privée de 1 à
10~mégabits par seconde, et publie trois latences par route~: radio à radio, rack à rack, et le chiffre que garantit son accord de
niveau de service (SLA). De Markham à Mahwah, elles valent 1.88, 1.89 et \qty{1.90}{\milli\second}.

La page ne précise pas s'il s'agit de temps à l'aller ou d'allers-retours, et elle n'a pas besoin de le faire~: la géodésique répond. De Markham à Mahwah, la lumière dans
le vide a besoin de \qty{1.747}{\milli\second} à l'aller, si bien qu'un aller-retour ne peut pas durer moins de \qty{3.49}{\milli\second}, et \qty{1.89}{\milli\second}
ne peut donc être qu'un temps d'aller simple, 8~\% au-dessus du \omterm{def:nw:the-north-american-map:floor}{plancher de latence} dans l'air. Le chiffre rack à rack ajoute \qty{10}{\micro\second} à Mahwah et \qty{20}{\micro\second} à
Carteret par rapport à la mesure radio à radio~: les courtes liaisons à chaque extrémité, du toit ou de la tour du bâtiment au rack du client.

\begin{remark}[Le reste du continent]\label{rem:nw:the-north-american-map:rest}
La méthode de ce chapitre s'étend à tout site dont le bâtiment peut être identifié~: les bourses d'options, les autres plateformes de contrats à terme, les marchés
du Mexique. \texttt{firm.venuesites} (chapitre~11) contient le registre qui relie chaque plateforme à son site~; les chapitres~22 à~26 le remplissent pour les plateformes
dont chaque classe d'actifs a besoin. Ce qui ne s'étend pas, c'est la certitude~: pour certains bâtiments, un dépôt ou la page d'un exploitant donne l'adresse postale, pour
d'autres seul un annuaire la donne, et la table des sites doit préciser lesquels.
\end{remark}

\section{Distances, temps de propagation et facteurs de tracé}

\begin{definition}[Plancher de latence, indice de groupe]\label{def:nw:the-north-american-map:floor}
Le \emph{plancher de latence}\index{plancher de latence} entre deux sites est le temps d'aller que met la lumière le long de la géodésique qui les relie, $d/c_0$ dans
le vide, avec $c_0 = \qty{299792458}{\metre\per\second}$~; dans un milieu, il vaut $n_g d / c_0$. L'\emph{indice de groupe}\index{indice de groupe} $n_g$ d'un
milieu est le rapport entre $c_0$ et la vitesse à laquelle une impulsion, et donc l'information, se propage dans ce milieu~: environ 1.462 dans une fibre monomode standard
à \qty{1550}{\nano\metre}, environ 1.0003 dans l'air près du sol.
\end{definition}

\begin{definition}[Facteur de tracé]\label{def:nw:the-north-american-map:factor}
Le \emph{facteur de tracé}\index{facteur de tracé} d'une route est sa latence publiée ou mesurée à l'aller, divisée par le \omterm{def:nw:the-north-american-map:floor}{plancher de latence} dans le
milieu propre à la route, entre ses deux extrémités. Un facteur de tracé de 1 désigne une route suivant la géodésique, sans équipement~; tout le reste est au-dessus de 1.
\end{definition}

\begin{proposition}[Ce que valent les planchers]\label{prop:nw:the-north-american-map:floors}
Par \qty{1000}{\kilo\metre} de géodésique, le \omterm{def:nw:the-north-american-map:floor}{plancher de latence} vaut \qty{3.336}{\milli\second} dans le vide et dans l'air, à
\qty{1}{\micro\second} près, et \qty{4.877}{\milli\second} dans la fibre~: sur un même tracé, le verre est 46~\% plus lent que l'air.
Une route de facteur $\rho$ dans un milieu d'\omterm{def:nw:the-north-american-map:floor}{indice de groupe} $n_g$ consacre $(\rho - 1)$ fois son plancher à tout ce qui n'est pas la géodésique, soit l'équivalent de
$(\rho - 1)\,d$ kilomètres de trajet supplémentaire dans ce milieu.
\end{proposition}

\begin{proof}
$10^3/c_0$ vaut \qty{3.3356}{\micro\second} par kilomètre~; multiplié par $1.0003$, il croît de \qty{1.0}{\micro\second} par \qty{1000}{\kilo\metre}~; multiplié par
1.4620, il vaut \qty{4.877}{\micro\second\per\kilo\metre}, soit 46.2~\% de plus. Un temps $\rho\, n_g d / c_0$ dépasse le plancher de $(\rho-1)\,n_g d/c_0$,
c'est le temps que met le même milieu pour $(\rho-1)\,d$ kilomètres de plus.
\end{proof}

L'indice propre de l'air compte à peine~: l'équation de Ciddor, telle que la tabule le NIST, donne 1.00027 pour l'air sec, à \qty{20}{\celsius} et
à la pression au niveau de la mer, et 1.00030 pour un air chaud, humide et comprimé, un écart qui déplace le plancher Aurora--Carteret de \qty{0.1}{\micro\second}.
Ce qui compte, c'est le verre face à l'air, soit \qty{1.82}{\milli\second} sur cette route, et l'écart dont la route s'écarte de la géodésique.

\omcode{code/firm/geomap/firm_geomap.py}{83}{89}{Planchers de latence et facteurs de tracé~: l'indice du milieu multiplié par la géodésique, divisé par la vitesse
de la lumière, et une latence publiée rapportée au plancher.}\label{lst:nw:the-north-american-map:floor}

Le \Cref{tab:nw:the-north-american-map:routes} applique les définitions aux routes publiées dans le registre. La fibre de 2010 de Laughlin, Aguirre et Grundfest,
entre 350 Cermak et Carteret, affichait une latence publiée de \qty{6.65}{\milli\second}, puis, plus tard, de \qty{6.55}{\milli\second}~: un \omterm{def:nw:the-north-american-map:factor}{facteur de tracé} de 1.21,
puis de 1.19, par rapport au plancher de la fibre, soit \qty{1.1}{\milli\second} de temps en trop, l'équivalent de 234 kilomètres de verre. Une partie tient au tracé (un
câble suit les routes et les emprises), une autre à l'équipement (amplificateurs, régénérateurs, commutateurs à chaque extrémité). Les mêmes auteurs ont évalué les
réseaux de liaisons hertziennes micro-ondes autorisés en 2011 et 2012 entre 4.2 et \qty{5.2}{\milli\second}~; le réseau qui atteignit Carteret en \qty{3.982}{\milli\second}
en 2016 avait un \omterm{def:nw:the-north-american-map:factor}{facteur de tracé} de 1.011, soit \qty{42}{\micro\second} au-dessus du plancher dans l'air, ou environ 12.5~kilomètres de détour et d'équipement sur
1\,181. Le chapitre~13 démonte les routes en fibre et le chapitre~14 les tours.

\begin{table}[ht]
\centering\small
\begin{tabular}{lllrrrr}
\toprule
Route & Origine & Destination & Publiée & Plancher & Facteur de & Écart \\
 & & & (ms) & (ms) & tracé & (\si{\micro\second}) \\
\midrule
Micro-ondes, 2016 & Aurora & Carteret & 3.982 & 3.940 & 1.011 & 42 \\
Micro-ondes, 2016 & Aurora & Mahwah & 3.986 & 3.934 & 1.013 & 52 \\
Sans fil, rack à rack & Markham & Mahwah & 1.890 & 1.748 & 1.082 & 142 \\
Sans fil, rack à rack & Markham & NY4 & 1.990 & 1.836 & 1.084 & 154 \\
Sans fil, rack à rack & Markham & Carteret & 2.050 & 1.845 & 1.111 & 205 \\
Fibre, 2010 & 350 Cermak & Carteret & 6.650 & 5.507 & 1.208 & 1\,143 \\
Fibre, plus tard & 350 Cermak & Carteret & 6.550 & 5.507 & 1.189 & 1\,043 \\
\bottomrule
\end{tabular}
\caption{Latences publiées à l'aller, comparées au plancher dans le milieu de chaque route (air pour les micro-ondes et le sans-fil, fibre pour la fibre). Sources~:
le communiqué de presse 2016 de Quincy Data pour les routes hertziennes, la page d'ICE sur le sans-fil à Toronto pour les routes de Markham, Laughlin, Aguirre et Grundfest
(2014) pour la fibre. Données~: \texttt{nw\_map.route\_rows()}.}\label{tab:nw:the-north-american-map:routes}
\end{table}

\omcode{code/networks/10-the-north-american-map/python/nw_map.py}{63}{75}{Chaque route publiée comparée à son plancher~: le facteur de tracé, le temps
en excès et le trajet supplémentaire dans le milieu propre à la route que cet excès représente.}\label{lst:nw:the-north-american-map:routes}

\begin{omfigure}
\begin{tikzpicture}
\begin{axis}[width=11.5cm,height=6.8cm,xmin=0,xmax=1250,ymin=0,ymax=7.2,xlabel={distance géodésique (km)},
  ylabel={latence à l'aller (ms)},tick label style={font=\footnotesize,/pgf/number format/1000 sep={\,}},label style={font=\small},grid=major,grid style={gray!25},
  legend pos=south east,legend style={font=\footnotesize},table/col sep=comma]
\addplot[omMeth,thick,dashed,no markers] table[x=km,y=fibre_ms]{figdata/networks/10-the-north-american-map/floors.csv};
\addplot[black,thick,no markers] table[x=km,y=vacuum_ms]{figdata/networks/10-the-north-american-map/floors.csv};
\addplot[only marks,mark=*,mark size=2.2pt,omThm] table[x=km,y=published_ms]{figdata/networks/10-the-north-american-map/routes_air.csv};
\addplot[only marks,mark=square*,mark size=2.2pt,omDef] table[x=km,y=published_ms]{figdata/networks/10-the-north-american-map/routes_fibre.csv};
\legend{plancher dans la fibre,plancher dans le vide,publié~: air,publié~: fibre}
\end{axis}
\end{tikzpicture}

\omcaption{Latences publiées à l'aller des routes en fonction de la \omterm{def:nw:the-north-american-map:geodesic}{distance géodésique} entre leurs extrémités, avec les planchers dans le vide (le plancher
dans l'air ne se distingue pas à cette échelle) et dans la fibre. Les routes sans fil se situent juste au-dessus de la droite du vide~; la fibre de 2010
se situe au-dessus même du plancher de la fibre. Données~: \texttt{fig\_map.py}, routes présentées dans le \cref{tab:nw:the-north-american-map:routes}.}\label{fig:nw:the-north-american-map:routes}
\end{omfigure}

\begin{remark}[Pourquoi l'ellipsoïde]\label{rem:nw:the-north-american-map:sphere}
La formule sphérique (haversine), avec le rayon moyen de la Terre, place Aurora et Carteret à \qty{2.97}{\kilo\metre} plus près que ne le fait l'ellipsoïde, une
erreur de 0.25~\% et \qty{9.9}{\micro\second} de lumière dans l'air~: plus d'un quart de l'écart entre la route hertzienne de 2016 et son plancher.
Entre Mahwah et Carteret, la sphère se trompe de 0.12~\% dans l'autre sens. Des \omterm{def:nw:the-north-american-map:factor}{facteurs de tracé} de 1.01 ne se mesurent pas avec une règle fausse de
0.25~\%, raison pour laquelle \texttt{firm.geomap} calcule la géodésique de Vincenty sur le WGS-84 et la teste contre l'ensemble publié par Karney de géodésiques.
\end{remark}

\begin{omfigure}
\begin{tikzpicture}[font=\footnotesize,>=Stealth]
\draw[thick,gray!60,dashed] (0,0) -- (11,0);
\node[below,gray] at (5.5,-0.1) {géodésique~: le plancher de latence};
\draw[thick,omDef] (0,0) -- (0.6,0.9) -- (2.4,1.2) -- (4.1,0.95) -- (5.9,1.35) -- (7.7,1.15) -- (9.3,1.0) -- (10.4,0.9) -- (11,0);
\foreach \x/\y in {0.6/0.9,2.4/1.2,4.1/0.95,5.9/1.35,7.7/1.15,9.3/1.0,10.4/0.9}{\draw[omThm,fill=omThm!30] (\x-0.09,\y-0.09) rectangle (\x+0.09,\y+0.09);}
\node[draw,fill=white,rounded corners=2pt,inner sep=2pt,align=center] at (0,-0.9) {rack dans\\ le centre de données A};
\node[draw,fill=white,rounded corners=2pt,inner sep=2pt,align=center] at (11,-0.9) {rack dans\\ le centre de données B};
\draw[<->,omMeth] (0.1,0.12) -- node[left,align=right,font=\scriptsize]{queue} (0.52,0.76);
\draw[<->,omMeth] (10.9,0.12) -- node[right,align=left,font=\scriptsize]{queue} (10.48,0.76);
\node[omThm,font=\scriptsize] at (5.9,1.8) {pylônes, amplificateurs ou régénérateurs};
\node[omDef,font=\scriptsize] at (5.5,0.45) {trajet~: détours autour du relief et des servitudes de passage};
\end{tikzpicture}

\omcaption{Ce que contient un \omterm{def:nw:the-north-american-map:factor}{facteur de tracé}, de façon schématique~: le détour du trajet par rapport à la géodésique, les équipements le long de celui-ci (pylônes, amplificateurs,
régénérateurs) et les queues à chaque extrémité, du toit du bâtiment ou de la \omterm{def:nw:colocation-products-and-how-they-are-sold:mmr}{salle d'interconnexion} jusqu'au rack du client. Un chiffre publié «~radio à
radio~» laisse les queues de côté~; «~rack à rack~» les inclut.}\label{fig:nw:the-north-american-map:parts}
\end{omfigure}

\section{Méthode~: une table de sites que l'on peut défendre}

\begin{method}[Construire la table des sites]\label{met:nw:the-north-american-map:table}
\begin{enumerate}
\item \textbf{Se procurer le bâtiment, non la ville.} Prendre le site dans les documents propres à la plateforme (un manuel de connectivité, un dépôt réglementaire, un rapport annuel)
  et l'adresse postale auprès de la plateforme ou de l'opérateur du centre de données. Un annuaire n'est qu'une piste, non une source~; s'il est la seule référence disponible, il faut signaler le site.
\item \textbf{Géocoder avec un géocodeur cité.} Consigner le géocodeur, la requête et la date~; conserver la latitude et la longitude à sept décimales
  (environ un centimètre), ce qui dépasse ce que justifie l'adresse, mais garantit la reproductibilité du calcul.
\item \textbf{Recouper de façon indépendante.} Comparer avec une position publiée, indépendante, du même bâtiment~; un écart de plus de quelques
  centaines de mètres signifie que l'une des deux est fausse.
\item \textbf{Calculer, ne jamais saisir à la main.} Les distances, les planchers et les coordonnées de carte proviennent de la table par l'intermédiaire du code~; les lignes de la table portent leurs identifiants du registre comptable
  afin que chaque nombre d'une carte puisse être rattaché à sa source.
\item \textbf{Dater la table.} Les plateformes déménagent (CME en 2010 et 2012, comme indiqué plus haut~; Euronext en 2022, au chapitre~11)~; il faut vérifier de nouveau la table avant
  chaque usage qui compte.
\end{enumerate}
\end{method}

Combien coûte une erreur sur un site~? Un kilomètre représente \qty{3.3}{\micro\second} de lumière dans le vide, ce qui correspond à l'ordre de grandeur des écarts que se disputent les
routes les plus rapides. La table des sites de ce chapitre satisfait à la troisième étape de la méthode~: Laughlin, Aguirre et Grundfest situent Aurora vers
41.80\textdegree{}~N, 88.24\textdegree{}~W et Carteret à 40.58\textdegree{}~N, 74.25\textdegree{}~W, ce qui concorde avec les coordonnées géocodées à moins
d'un centième de degré près (moins d'un kilomètre). Leur distance, \qty{1179}{\kilo\metre}, se situe entre les distances sphérique (\qty{1177.2}{\kilo\metre}) et ellipsoïdale
(\qty{1180.1}{\kilo\metre}) entre leurs positions arrondies, et à moins de \qty{2}{\kilo\metre} de la nôtre~; l'article ne dit pas quelle figure de la
Terre il a utilisée, et son temps de lumière de \qty{3.93}{\milli\second} concorde avec le nôtre à \qty{6}{\micro\second} près. Markham est la ligne faible~: son
bâtiment est sourcé auprès de TMX et d'ICE, son adresse postale seulement auprès d'un annuaire, et une erreur de quelques kilomètres sur Markham déplacerait ses planchers
d'environ \qty{10}{\micro\second}, que les \omterm{def:nw:the-north-american-map:factor}{facteurs de tracé} du \cref{tab:nw:the-north-american-map:routes} absorberaient sans que personne ne le remarque.

\omcode{code/firm/geomap/firm_geomap.py}{43}{73}{Inversion de Vincenty sur l'ellipsoïde WGS-84, avec une garde de convergence~: l'itération sur la
différence de longitude sur la sphère auxiliaire, puis la série pour la distance.}\label{lst:nw:the-north-american-map:vincenty}

\section{Tutoriel~: la carte à partir des coordonnées calculées}\label{tut:nw:the-north-american-map:tut}

\begin{tutorial}
\textbf{Objectif.} Charger la table des sites, calculer chaque distance et chaque plancher, comparer les routes publiées à leurs planchers et tracer les cartes.
\textbf{État final~:} \cref{fig:nw:the-north-american-map:nj,fig:nw:the-north-american-map:na,fig:nw:the-north-american-map:routes} et
\cref{tab:nw:the-north-american-map:pairs,tab:nw:the-north-american-map:routes}.
\begin{tutsteps}
\item \textbf{La table.} \texttt{data/networks/sites.csv} contient sept sites, chacun avec les lignes du registre qui sourcent son bâtiment, son adresse et
  ses coordonnées~; \texttt{firm\_geomap.load\_sites} la lit.
\item \textbf{L'étalon.} Lancer les tests de \texttt{firm.geomap}~: la distance de Vincenty
  (\cref{lst:nw:the-north-american-map:vincenty}) contre 249 géodésiques de l'ensemble de test de Karney, en Python et en C++20, à mieux qu'un dixième
  de millimètre près.
\item \textbf{Les planchers.} \texttt{nw\_map.pair\_rows()} calcule le \cref{tab:nw:the-north-american-map:pairs}~; \texttt{route\_rows()}
  (\cref{lst:nw:the-north-american-map:routes}) calcule le \cref{tab:nw:the-north-american-map:routes}.
\item \textbf{Les cartes.} \texttt{python fig\_map.py} projette les sites de chaque région et écrit les fichiers de sites, de liaisons et d'étiquettes que lisent les cartes TikZ~;
  aucune coordonnée n'est saisie dans la figure.
\end{tutsteps}
\textbf{Pour aller plus loin.} Ajouter les points de présence du site secondaire de Cboe, ou une autre plateforme d'après son propre manuel, avec sa ligne du registre, et
recalculer~; remplacer l'indice de la fibre par celui d'une fibre à cœur creux (chapitre~13) et observer ce qui arrive au plancher de Chicago.
\end{tutorial}

\section{Construction~: la carte géodésique}\label{bld:nw:the-north-american-map:geomap}

\begin{build}
\bfield{Rôle} Des sites tirés de données sourcées, des distances sur l'ellipsoïde, des \omterm{def:nw:the-north-american-map:floor}{planchers de latence} et des \omterm{def:nw:the-north-american-map:factor}{facteurs de tracé}, et des coordonnées pour les cartes~: la base de
toutes les cartes et routes de la partie~III, ainsi que du plan du chapitre~29.
\bfield{Interface} \texttt{firm\_geomap}~: \texttt{Site(id, name, operator, lat, lon, source, venues)}, \texttt{geodesic\_m},
\texttt{haversine\_m}, \texttt{floor\_us(d\_m, medium, n\_g)}, \texttt{route\_factor(published\_us, d\_m, medium)}, \texttt{matrix},
\texttt{project(sites, lat0, lon0)}, \texttt{load\_sites}, \texttt{write\_sites}~; une jumelle C++20 de la géodésique dans \texttt{cpp/firm\_geomap.hpp}.
\bfield{Règles} Chaque site nomme ses lignes de source~; la géodésique lève une erreur plutôt que de renvoyer une distance fausse lorsque l'itération échoue (points
presque antipodaux)~; les planchers sont donnés à l'aller~; une projection ne sert qu'au tracé, jamais à calculer une distance.
\bfield{Tests d'acceptation} \texttt{code/firm/geomap/tests/}~: 249 géodésiques de l'ensemble de test de Karney, sous licence CC0, à 0.1~mm près en Python et en
C++20, les planchers vérifiés à la main, la sphère comparée à l'ellipsoïde, et la présence des sources de la table des sites.
\bfield{Pour aller plus loin} L'algorithme de Karney, qui converge partout~; un chemin en grand cercle sous forme de liste de points, pour tracer une route sur une carte.
\end{build}

\begin{omsources}
\item G.~Laughlin, A.~Aguirre et J.~Grundfest, «~Information transmission between financial markets in Chicago and New York~», \emph{Financial
Review} 49(2) (2014).
\item T.~Vincenty, «~Direct and inverse solutions of geodesics on the ellipsoid with application of nested equations~», \emph{Survey Review}
23(176) (1975)~; C.~F.~F.~Karney, «~Algorithms for geodesics~», \emph{Journal of Geodesy} 87 (2013), et l'ensemble de test pour les géodésiques (Zenodo, CC0).
\item ICE, \emph{Colocation Technical Specifications} (Mahwah) et la page wireless de Toronto~; Cboe, \emph{Titanium U.S. Equities/Options
Connectivity Manual}~; CME Group, Form 10-K 2024~; SEC Release 34-101267 (Nasdaq)~; Quincy Data, publication du 13 mai 2016.
\item BIPM, la définition du mètre dans le SI~; NIST, indice de réfraction de l'air (équation de Ciddor)~; coordonnées des sites tirées d'OpenStreetMap Nominatim.
\end{omsources}

\section{Exercices}

\begin{exercise}[$\star$]\label{exo:nw:the-north-american-map:1}
Calculer le temps que met la lumière pour parcourir \qty{1000}{\kilo\metre} dans le vide et dans une fibre standard, puis la différence sur le trajet Aurora--Carteret.
\end{exercise}

\begin{exercise}[$\star$]\label{exo:nw:the-north-american-map:2}
Un fournisseur publie une latence de \qty{2.05}{\milli\second} entre Markham et Carteret sans dire s'il s'agit d'un trajet simple ou d'un aller-retour. Montrer de quel cas il relève.
\end{exercise}

\begin{exercise}[$\star$]\label{exo:nw:the-north-american-map:3}
La formule sphérique place Aurora et Carteret \qty{2.97}{\kilo\metre} trop près l'une de l'autre. Calculer combien de microsecondes cela représente dans l'air, et expliquer pourquoi cela compte.
\end{exercise}

\begin{exercise}[$\star\star$]\label{exo:nw:the-north-american-map:4}
Le tracé hertzien micro-ondes de 2016 entre Aurora et Mahwah a été publié à \qty{3.986}{\milli\second}. Calculer son \omterm{def:nw:the-north-american-map:factor}{facteur de tracé} et le trajet supplémentaire
qu'il représente dans l'air.
\end{exercise}

\begin{exercise}[$\star\star$]\label{exo:nw:the-north-american-map:5}
Un routeur à Carteret envoie une partie d'un ordre vers Secaucus NY5 et une autre vers Mahwah par une fibre en ligne droite. Calculer de combien le routage synchronisé
doit retarder la partie destinée à NY5, et expliquer pourquoi la valeur réelle est différente.
\end{exercise}

\begin{exercise}[$\star\star$]\label{exo:nw:the-north-american-map:6}
Déterminer quelle ligne de la table des sites est la moins fiable, et ce qu'une erreur de 3~kilomètres dans cette ligne ferait aux planchers qui l'utilisent.
\end{exercise}

\begin{exercise}[$\star\star\star$]\label{exo:nw:the-north-american-map:7}
\emph{Programmation.} Arrondir les coordonnées d'Aurora et de Carteret à deux décimales, comme les impriment Laughlin, Aguirre et Grundfest, puis calculer la géodésique
avec \texttt{firm.geomap}, puis mesurer l'écart avec celle de leur article et avec celle obtenue sans arrondi.
\end{exercise}

\begin{exercise}[$\star\star\star$]\label{exo:nw:the-north-american-map:8}
\emph{Trouver la faille.} «~Notre tracé de fibre entre Chicago et le New Jersey a un \omterm{def:nw:the-north-american-map:factor}{facteur de tracé} de 1.19, donc un réseau micro-ondes sur le même trajet serait
19~\% plus rapide.~»
\end{exercise}

\section{Problème~: de Carteret à Aurora}

\begin{problem}[{Problème du week-end --- les planchers entre les centres de données du Nasdaq et du CME, et un tracé en fibre confronté à ces planchers}]\label{pb:nw:the-north-american-map:1}
Une firme négocie des contrats à terme sur indices à Aurora contre des actions à Carteret et veut savoir à quelle vitesse la liaison entre les deux pourrait au mieux fonctionner, et à quelle distance
de cette limite les tracés publiés s'en sont écartés.

\medskip\noindent\textbf{Partie I --- La distance.}
\begin{enumerate}
\item D'après la table des sites, quelle est la \omterm{def:nw:the-north-american-map:geodesic}{distance géodésique} entre Aurora et Carteret~?
\item Que donne la formule sphérique, et quelle est son erreur en kilomètres et en pourcentage~?
\item De combien la corde droite à travers la Terre est-elle plus courte~?
\item Pourquoi la géodésique, et non la corde, est-elle la bonne mesure pour un tracé de trading~?
\end{enumerate}

\medskip\noindent\textbf{Partie II --- Les planchers.}
\begin{enumerate}[resume]
\item Quel est le \omterm{def:nw:the-north-american-map:floor}{plancher de latence} à l'aller dans le vide~?
\item Dans l'air, en prenant $n = 1.0003$, quelle importance a le choix entre 1.00027 et 1.0003~?
\item Dans une fibre standard, le long de la géodésique~?
\item Quel est le plancher d'aller-retour dans la fibre, et dans le vide~?
\end{enumerate}

\medskip\noindent\textbf{Partie III --- Un tracé en fibre face à son plancher.}
\begin{enumerate}[resume]
\item La fibre de 2010 entre 350 Cermak et Carteret a été cotée à \qty{6.65}{\milli\second} à l'aller. Quel est son plancher, et quel est son \omterm{def:nw:the-north-american-map:factor}{facteur de tracé}~?
\item Quel est ce temps en excès, et à combien de kilomètres de fibre cela équivaut-il~?
\item Depuis Aurora, une firme rejoignait cette fibre par le réseau métropolitain Cermak--Aurora. Si le tronçon métropolitain avait le même \omterm{def:nw:the-north-american-map:factor}{facteur de tracé}, quel serait le temps
  d'Aurora à Carteret, et son \omterm{def:nw:the-north-american-map:factor}{facteur de tracé} par rapport au plancher Aurora--Carteret dans la fibre~?
\item Le tracé a ensuite été coté à \qty{6.55}{\milli\second}. Qu'est-ce qui a changé dans le \omterm{def:nw:the-north-american-map:factor}{facteur de tracé}~?
\end{enumerate}

\medskip\noindent\textbf{Partie IV --- Le verdict.}
\begin{enumerate}[resume]
\item Énoncer le \emph{résultat nommé}~: les trois planchers entre Carteret et Aurora et le \omterm{def:nw:the-north-american-map:factor}{facteur de tracé} de la fibre de 2010.
\item Le tracé hertzien micro-ondes de 2016 a atteint Carteret en \qty{3.982}{\milli\second}. De combien a-t-il battu le plancher de la fibre~?
\item Quel écart reste-t-il entre lui et le plancher dans l'air~?
\item Un tracé en fibre pourrait-il jamais le battre~? Quel type de fibre faudrait-il pour y parvenir (chapitre~13)~?
\item Quelle erreur sur un site déplacerait le résultat nommé de plus de \qty{1}{\micro\second}~?
\item Que fait le déplacement du moteur d'appariement du CME de Cermak à Aurora à la valeur de la fibre ancrée à Cermak~?
\item Pourquoi les firmes publient-elles et achètent-elles des latences de baie à baie plutôt que des planchers~?
\item En une phrase~: que dit la géodésique à une firme que la latence d'un fournisseur ne dit pas~?
\end{enumerate}
\end{problem}

\section{Questions d'entretien}

\begin{interviewq}[$\star$ \iqroles{développeur, trader}]\label{iq:nw:the-north-american-map:1}
Où se trouvent les principaux moteurs d'appariement américains des actions, des options et des contrats à terme, et à quelle distance environ les uns des autres~?
\end{interviewq}

\begin{interviewq}[$\star$ \iqroles{développeur}]\label{iq:nw:the-north-american-map:2}
Quelle est la latence minimale à l'aller entre Chicago et le New Jersey, dans le vide et dans la fibre~? Montrez le calcul.
\end{interviewq}

\begin{interviewq}[$\star\star$ \iqroles{développeur, chercheur}]\label{iq:nw:the-north-american-map:3}
Pourquoi une liaison hertzienne micro-ondes est-elle plus rapide qu'une fibre sur le même trajet, et pourquoi utilise-t-on encore la fibre~?
\end{interviewq}

\begin{interviewq}[$\star\star$ \iqroles{développeur}]\label{iq:nw:the-north-american-map:4}
Comment calculeriez-vous la distance entre deux centres de données, et avec quelle précision doit-elle être~?
\end{interviewq}

\begin{interviewq}[$\star\star$ \iqroles{trader, chercheur}]\label{iq:nw:the-north-american-map:5}
Pourquoi un trader à Carteret voit-il les cotations du Nasdaq avant que le flux consolidé ne les affiche, et quelle part de l'écart relève de la physique~?
\end{interviewq}

\begin{interviewq}[$\star\star\star$ \iqroles{développeur, trader}]\label{iq:nw:the-north-american-map:6}
Un fournisseur propose une liaison Chicago--New Jersey à \qty{3.99}{\milli\second}. Comment jugeriez-vous cette affirmation avant d'acheter~?
\end{interviewq}
